% ══════════════════════════════════════════════════════════════════════════════
%  EXPERIMENT 5 — MEASURING g                         (Lesson 13, T3 Week 5)
%  From Lesson 13 of the A4 Forces Lesson Booklet; answer keys 13.* kept.
% ══════════════════════════════════════════════════════════════════════════════
\experimentdiv{Measuring \texorpdfstring{$g$}{g}}{A motion sensor \textbullet{} Velocity--time graphs \textbullet{} Anomalies and percentage error}{13}\label{sec:5}
\lessoncodes{VC2S8U14 \textbullet{} VC2S8I01, I02, I03, I04, I05, I06}

\vocabhead
\begin{ctab}
\begin{tabular}{|SW{4.2cm}|W{11.6cm}|}
  \hline\hd & \thead{Word} & \thead{Definition (copy from the slide)} \\ \hline
  & \scaf{\textbf{Motion sensor}} & \scaf{A device that sends out pulses of sound and times their echoes, to measure how far away an object is many times a second.} \\ \hline
  & \textbf{Velocity} & \cell{13.v2} \\ \hline
  & \textbf{Acceleration} & \cell{13.v6} \\ \hline
  & \textbf{Accepted value} & \cell{13.v3} \\ \hline
  & \textbf{Percentage error} & \cell{13.v4} \\ \hline
\end{tabular}
\end{ctab}

\begin{lipanel}
\inquiryli{Collect data with a digital sensor and judge how accurate it is.}
\understandli{Test whether the mass of a falling object changes its acceleration.}
\sctitle
\begin{itemize}[leftmargin=1.4em, itemsep=0pt, topsep=1pt]
  \item I can read starting velocity, final velocity and time of fall off a graph.
  \item I can calculate an acceleration and a percentage error.
  \item I can say whether my results support Galileo's idea, using numbers.
\end{itemize}
\end{lipanel}

\begin{skillsfocus}
\focus{QP1}{Write a hypothesis in the faculty frame.}
\focus{PC8}{Set up and read the motion sensor and its software.}
\focus{PA1}{Record every trial in a THUD table.}
\focus{PA3}{Calculate accelerations, means and percentage errors.}
\focus{PA7}{Spot an anomaly and leave it out of the mean.}
\focus{EV2}{Compare your value with the accepted value of $g$.}
\focus{EV3}{Name two errors and how to reduce them.}
\end{skillsfocus}

% ─── READING ─────────────────────────────────────────────────────────────────
\exppart{Before the lab}{Teal}

\begin{readingbox}{Reading --- The Hammer and the Feather}
A ball dropped from head height hits the floor in about half a second --- too quick to time with a stopwatch. For almost two thousand years, people trusted the Greek thinker Aristotle, who taught that heavy objects fall faster than light ones.

Galileo disagreed. A famous story says he dropped two balls of different mass from the Leaning Tower of Pisa, and they landed together. Without air in the way, heavy and light objects speed up at exactly the same rate. In 1971 the astronaut David Scott tested this on the Moon, where there is no air: he dropped a hammer and a feather, and they hit the ground at the same moment.

A motion sensor can measure a fall for you. Like a bat hunting in the dark, it sends out short pulses of sound, too high for us to hear, and times how long each echo takes to come back. It does this 50 times a second, so it knows how fast a falling object is moving at every moment.

In this experiment you will drop three weights of different mass under a motion sensor, find each one's acceleration from its graph, and test Galileo's idea. You will also compare your results with the accepted value of $g$, 9.8\,m/s$^2$.
\end{readingbox}

\rcreset
\rcqk{13.c1}{Why is a ball dropped from head height hard to time with a stopwatch?}{2.2cm}
\rcqk{13.c2}{What did Aristotle teach about falling objects, and what did Galileo say instead?}{2.4cm}
\rcqk{13.c3}{What happened when David Scott dropped a hammer and a feather on the Moon? Why would the same test not work on Earth?}{2.4cm}
\rcqk{13.c4}{Explain how a motion sensor works out how far away an object is.}{2.2cm}

% ─── TECHNIQUE ────────────────────────────────────────────────────────────────
\exppart{Technique spotlight}{Blue}

\begin{technique}{Acceleration from a velocity--time graph}{PC8}
While the object falls, its speed goes up steadily, so the line on the graph is straight. Read three things off it: the \textbf{starting velocity} (the moment it is let go), the \textbf{final velocity} (just before it lands) and the \textbf{time of fall} between them. Then ask the software for the \emph{slope} of the straight part as a check.
\end{technique}

\begin{center}
\diagVelocityTime{parts}
\expfigcap{What one drop looks like on the velocity--time graph.}
\end{center}

\begin{workedex}{Acceleration from a graph}
A falling 100\,g weight: starting velocity 0.2\,m/s, final velocity 4.1\,m/s, time of fall 0.40\,s; the software's slope is 9.7\,m/s$^2$.

\smallskip
\textbf{Change in speed:}\quad $4.1 - 0.2 = 3.9$\,m/s \qquad \textbf{Acceleration:}\quad $3.9 \div 0.40 = \mathbf{9.75\ \text{m/s}^2}$\\
\textbf{Check:}\quad the slope, 9.7\,m/s$^2$, is very close, so the graph was read well.
\end{workedex}

\begin{wedo}{Acceleration from a graph}
Another drop shows a starting velocity of 0.1\,m/s, a final velocity of 4.5\,m/s and a time of fall of 0.45\,s.
\pt{a}Find the change in speed.
\wbox{13.wd1a}{1.8cm}
\pt{b}Find the acceleration, to one decimal place.
\wbox{13.wd1b}{1.8cm}
\end{wedo}

\begin{infobox}{Percentage error}
\[ \text{percentage error} = \frac{\text{your value} - \text{accepted value}}{\text{accepted value}} \times 100\% \]
Ignore any minus sign. For $g$, the accepted value is 9.8\,m/s$^2$. A smaller percentage error means a more accurate result.
\end{infobox}

\begin{workedex}{Percentage error}
A group's mean value of $g$ is 9.2\,m/s$^2$.

\smallskip
\textbf{Difference:}\quad $9.8 - 9.2 = 0.6$ \qquad \textbf{Percentage error:}\quad $0.6 \div 9.8 \times 100\% \approx \mathbf{6\%}$
\end{workedex}

\begin{wedo}{Percentage error}
Another group's mean value of $g$ is 10.3\,m/s$^2$.
\pt{a}Find the difference from the accepted value of 9.8\,m/s$^2$.
\wbox{13.wd2a}{1.6cm}
\pt{b}Find the percentage error.
\wbox{13.wd2b}{1.8cm}
\end{wedo}

% ─── PLAN, SAFETY AND METHOD ─────────────────────────────────────────────────
\exppart{Plan, safety and method}{Teal}

\textbf{Aim:} To investigate the effect of changing the mass of a falling object on its acceleration.

\task Complete the variables table. \skr{PC2}
\begin{ctab}
\begin{tabular}{|SW{4.4cm}|W{10.2cm}|}
  \hline\hd & \thead{Variable} & \thead{In this experiment} \\ \hline
  & \scaf{Independent} & \scaf{The mass of the weight: 50, 100 and 200\,g} \\ \hline
  & Dependent & \cell{13.t1-r2} \\ \hline
  & Controlled & \cell{13.t1-r3} \\ \hline
  & Controlled & \cell{13.t1-r4} \\ \hline
\end{tabular}
\end{ctab}

\needspace{8\baselineskip}
\task Write your hypothesis. \skr{QP1}

\scaf{If the mass of the falling weight is increased, its acceleration will \dots\ because \dots}
\abox{13.pred}{2.5cm}

\subsubsection*{Equipment}
Motion sensor and data logger (or a computer with the logging software) \textbullet{} stand, rod and right-angle clamp \textbullet{} 50\,g, 100\,g and 200\,g weights \textbullet{} cushion or no-bounce pad \textbullet{} metre rule

\begin{center}
\diagMotionSensorSetup
\expfigcap{The sensor points straight down at a cushion, so the weight does not bounce back up towards it.}
\end{center}

\needspace{12\baselineskip}
\task Complete the risk assessment. \skr{PC3}
\begin{ctab}
\begin{tabular}{|SW{4.0cm}|W{5.2cm}|W{5.4cm}|}
  \hline\hd & \thead{Risk (what could go wrong)} & \thead{Prevention} & \thead{Responding (if it happens)} \\ \hline
  & \scaf{The stand tips over} & \scaf{Stand the base well onto the bench} & \scaf{Step back; tell the teacher; check for injuries} \\ \hline
  & A falling weight hits a foot or hand & & \\ \hline
  & Someone climbs to reach the sensor & & \\ \hline
\end{tabular}
\end{ctab}

\needspace{18\baselineskip}
\subsubsection*{Method}
\begin{enumerate}[itemsep=2pt]
  \item Clamp the motion sensor to the stand at the edge of the bench, as high as it will go, pointing straight down. Put the cushion on the floor directly underneath it.
  \item Connect the sensor to the data logger. Set the sensor to long range, set the logger to 50 readings per second, and show a graph of velocity against time.
  \item Hold the 50\,g weight flat, directly under the sensor, with its bottom 40\,cm below the sensor. Keep your hand out of the way of the sound pulses.
  \item Start recording, then let go --- do not push it. Stop recording once it has landed.
  \item From the graph, read the starting velocity, the final velocity and the time of fall, and ask the software for the slope. Record all four in Table 5.1.
  \item Repeat steps 3--5 until you have three trials, then do the same for the 100\,g and 200\,g weights.
\end{enumerate}

% ─── RESULTS ─────────────────────────────────────────────────────────────────
\exppart{Results}{Amber}

\subsubsection*{Table 5.1: Velocity, time of fall and acceleration for each trial \skr{PA1}}
\begin{ctab}
\begin{tabular}{|SZ{2.3cm}|Z{2.0cm}|Z{2.0cm}|Z{1.9cm}|Z{3.6cm}|Z{2.3cm}|}
  \hline\hd & \thead{Weight, trial} & \thead{Starting velocity (m/s)} & \thead{Final velocity (m/s)} & \thead{Time of fall (s)} & \thead{Acceleration (m/s$^2$)} & \thead{Slope from the software (m/s$^2$)} \\ \hline
  & \scaf{e.g.\ 50\,g, 1} & \scaf{0.1} & \scaf{4.3} & \scaf{0.45} & \scaf{$(4.3 - 0.1) \div 0.45 = 9.3$} & \scaf{9.4} \\ \hline
  & 50\,g, 1 & \cell{13.res-r2c2} & \cell{13.res-r2c3} & \cell{13.res-r2c4} & \cell{13.res-r2c5} & \cell{13.res-r2c6} \\ \hline
  & 50\,g, 2 & \cell{13.res-r3c2} & \cell{13.res-r3c3} & \cell{13.res-r3c4} & \cell{13.res-r3c5} & \cell{13.res-r3c6} \\ \hline
  & 50\,g, 3 & \cell{13.res-r4c2} & \cell{13.res-r4c3} & \cell{13.res-r4c4} & \cell{13.res-r4c5} & \cell{13.res-r4c6} \\ \hline
  & 100\,g, 1 & \cell{13.res-r5c2} & \cell{13.res-r5c3} & \cell{13.res-r5c4} & \cell{13.res-r5c5} & \cell{13.res-r5c6} \\ \hline
  & 100\,g, 2 & \cell{13.res-r6c2} & \cell{13.res-r6c3} & \cell{13.res-r6c4} & \cell{13.res-r6c5} & \cell{13.res-r6c6} \\ \hline
  & 100\,g, 3 & \cell{13.res-r7c2} & \cell{13.res-r7c3} & \cell{13.res-r7c4} & \cell{13.res-r7c5} & \cell{13.res-r7c6} \\ \hline
  & 200\,g, 1 & \cell{13.res-r8c2} & \cell{13.res-r8c3} & \cell{13.res-r8c4} & \cell{13.res-r8c5} & \cell{13.res-r8c6} \\ \hline
  & 200\,g, 2 & \cell{13.res-r9c2} & \cell{13.res-r9c3} & \cell{13.res-r9c4} & \cell{13.res-r9c5} & \cell{13.res-r9c6} \\ \hline
  & 200\,g, 3 & \cell{13.res-r10c2} & \cell{13.res-r10c3} & \cell{13.res-r10c4} & \cell{13.res-r10c5} & \cell{13.res-r10c6} \\ \hline
\end{tabular}
\end{ctab}

\needspace{12\baselineskip}
\subsubsection*{Table 5.2: Mean acceleration and percentage error for each weight}
\begin{ctab}
\begin{tabular}{|SZ{2.4cm}|Z{2.2cm}|Z{5.6cm}|Z{4.4cm}|}
  \hline\hd & \thead{Weight} & \thead{Mass (kg)} & \thead{Mean acceleration (m/s$^2$), leaving out any anomalies} & \thead{Percentage error (\%)} \\ \hline
  & \scaf{e.g.\ 50\,g} & \scaf{0.05} & \scaf{$(9.3 + 9.5 + 9.4) \div 3 = 9.4$} & \scaf{$(9.8 - 9.4) \div 9.8 \times 100 \approx 4$} \\ \hline
  & 50\,g & \cell{13.mean-r2c2} & \cell{13.mean-r2c3} & \cell{13.mean-r2c4} \\ \hline
  & 100\,g & \cell{13.mean-r3c2} & \cell{13.mean-r3c3} & \cell{13.mean-r3c4} \\ \hline
  & 200\,g & \cell{13.mean-r4c2} & \cell{13.mean-r4c3} & \cell{13.mean-r4c4} \\ \hline
\end{tabular}
\end{ctab}

% ─── ANALYSIS ────────────────────────────────────────────────────────────────
\exppart{Analysis and evaluation}{Coral}

\task Does the mass of a falling object change its acceleration? Use your mean accelerations. \skr{PA5}
\abox{13.t2a}{3.5cm}

\needspace{12\baselineskip}
\task \skr{PA7}
\pt{a}Did any trial look like an anomaly? Say which, and suggest what might have gone wrong.
\abox{13.t4a}{2.4cm}
\pt{b}Why is it better to use the mean of three trials than a single trial?
\abox{13.t4b}{2.4cm}

\needspace{8\baselineskip}
\task A falling weight speeds up from 0.3\,m/s to 3.9\,m/s in 0.37\,s. Find its acceleration. \skr{PA3}
\wbox{13.t5}{2.6cm}

\needspace{8\baselineskip}
\task A group's mean acceleration was 9.5\,m/s$^2$. Calculate their percentage error. \skr{EV2}
\wbox{13.t6}{2.6cm}

\needspace{10\baselineskip}
\task An apple and a bowling ball are dropped from the same height at the same moment. Which lands first? Use your results to explain.
\abox{13.t7}{3.5cm}

\needspace{14\baselineskip}
\task Identify two sources of error. For each, suggest a change to the method that would reduce it. \skr{EV3}
\pt{a}First source of error and how to reduce it:
\abox{13.t8a}{2.4cm}
\pt{b}Second source of error and how to reduce it:
\abox{13.t8b}{2.4cm}

\needspace{12\baselineskip}
\task Write a conclusion: state what you found, back it with numbers from Table 5.2, and say whether your hypothesis was supported.

\scaf{The mean accelerations of the 50\,g, 100\,g and 200\,g weights were \dots, \dots\ and \dots\,m/s$^2$, within \dots\,\% of 9.8\,m/s$^2$. This supports / does not support the hypothesis.}
\abox{13.concl}{5cm}

\begin{yr8box}
\begin{keepext}
\exttask The motion sensor takes 50 readings every second. How much time passes between readings, and about how many readings does it take during a 0.45\,s fall? Explain why a sensor that took only 5 readings a second would give a much less precise value of $g$.
\wbox{13.x9}{4.5cm}
\end{keepext}
\begin{keepext}
\exttask Air resistance pushes up on a falling object. The 50\,g weight often has a slightly smaller acceleration than the 200\,g weight. Use the idea of net force to explain why.
\abox{13.x10}{4.5cm}
\end{keepext}
\end{yr8box}

% ─── SKILL CHECK ─────────────────────────────────────────────────────────────
\exppart{Skill check}{NavyBlue}

\begin{skillcheck}
\scrow{QP1}{My hypothesis names the mass and the acceleration.}
\scrow{PC8}{I set up the sensor and read its graph and slope.}
\scrow{PA1}{My tables have titles, headings, units in the headings and matching decimals.}
\scrow{PA3}{I calculated accelerations, means and percentage errors.}
\scrow{PA7}{I spotted any anomaly and left it out of the mean.}
\scrow{EV2}{I compared my value of $g$ with 9.8\,m/s$^2$.}
\scrow{EV3}{I named two errors and ways to reduce them.}
\end{skillcheck}

\task Look back over all five experiments. Which skill have you improved most this unit? Give an example.
\abox{13.sc}{5cm}

\end{document}
